\documentclass[10pt]{article} % For LaTeX2e
\usepackage{tmlr}
% If accepted, instead use the following line for the camera-ready submission:
%\usepackage[accepted]{tmlr}
% To de-anonymize and remove mentions to TMLR (for example for posting to preprint servers), instead use the following:
%\usepackage[preprint]{tmlr}

% Optional math commands from https://github.com/goodfeli/dlbook_notation.
\input{math_commands.tex}

\usepackage{hyperref}
\usepackage{url}
\usepackage{amssymb}
\usepackage{amsmath}
\usepackage{algorithm}
\usepackage[noend]{algpseudocode}
 \usepackage{multirow}
 \usepackage{subfig}
\usepackage{graphicx}
\usepackage{lscape}
\usepackage{xcolor}


\title{Prediction Interval Estimation in Support Vector Machines}


% Authors must not appear in the submitted version. They should be hidden
% as long as the tmlr package is used without the [accepted] or [preprint] options.
% Non-anonymous submissions will be rejected without review.

\author{\name Pritam Anand \email pritam$\_$anand@daiict.ac.in \\
      \addr DA-IICT, Gandhinagar.
     }

% The \author macro works with any number of authors. Use \AND 
% to separate the names and addresses of multiple authors.

\newcommand{\fix}{\marginpar{FIX}}
\newcommand{\new}{\marginpar{NEW}}

\def\month{MM}  % Insert correct month for camera-ready version
\def\year{YYYY} % Insert correct year for camera-ready version
\def\openreview{\url{https://openreview.net/forum?id=XXXX}} % Insert correct link to OpenReview for camera-ready version


\begin{document}


\maketitle

\begin{abstract}
This paper explores Uncertainty Quantification (UQ) in SVM predictions, particularly for regression and forecasting tasks. SVM models obtain promising prediction, comparable to modern complex Neural Network architectures, particularly for small-scale tabular datasets with global optimal solution. Unlike the Neural Network, there are  only few literature which details about the UQ in SVM prediction.  We provide a comprehensive summary of existing Prediction Interval (PI) estimation and probabilistic forecasting methods developed in the SVM framework. We begin by outlining key properties of an ideal PI estimation model and systematically evaluating existing PI estimation and probabilistic forecasting methods in SVM against these criteria. After a thorough investigation, we find that none of the existing SVM PI models achieves a sparse solution, which has remained a key advantage of the standard SVM model developed for classification and regression tasks. To introduce sparsity in SVM model, we propose the Sparse Support Vector Quantile Regression (SSVQR) model, which constructs PIs and probabilistic forecasts by solving a pair of linear programs. Additionally, we develop a feature selection algorithm tailored for PI estimation using SSVQR. Our algorithm not only eliminates a significant number of features in high-dimensional data but also enhances the overall quality of the PI in case of high-dimensional data. Extensive experiments on artificial, real-world benchmark datasets empirically demonstrate and compare the different characteristics of both existing and proposed SVM-based PI estimation methods. Furthermore, we compare both, the existing and proposed SVM-based  PI estimation models, with modern deep learning models for probabilistic forecasting tasks on benchmark datasets. The results demonstrate that SVM-based models can achieve a comparable quality of probabilistic forecasts to that of complex deep learning architectures. 
\end{abstract}

\section{Introduction}
\label{sec:sample1}
Given the training set $T=\{(x_i,y_i): x_i \in \mathbf{R}^n, y_i \in \mathbf{R}, i =1,2,...m. \}$, sampled independently from the joint distribution of the random variables $(X,Y)$, the goal of the regression task is to estimate a function that predicts the target variable  $y$ based on the input variable $x$ well. However, in most of applications, the prediction of the regression model may not be perfectly accurate due to random relationship between $Y$ and $X$. For example, predicting the impact of a specific drug on a patient's heart rate based on their Body Mass Index (BMI) may  not be accurate and  may involve a significant degree of uncertainty. In such cases, quantifying these uncertainties is crucial for making effective decisions.

 The Prediction Interval (PI) estimation is most commonly used Uncertainty Quantification (UQ) technique in regression tasks. Given a high confidence $1-\alpha \in (0,1)$ and training set $T$, the PI  tube is defined as a pair of functions $(f_1,f_2)$. It is said to be well calibrated if it satisfies $P(f_1(X)\leq Y \leq f_2(X)|X) \geq 1-\alpha$. The objective of the PI models is to obtain a PI tube with the minimum possible width while ensuring the target calibration. Therefore, the performance of a PI estimation method is mainly evaluated using two criteria: Prediction Interval Coverage Probability (PICP), which computes the fraction of $y$ values within the PI tube, and Mean Prediction Interval Width (MPIW), quantifying the width of the PI tube.


The  PI estimation models requires to explore the different  characteristics of the conditional distribution $Y|X$, rather focusing only on $E(Y|X)$ as done in standard regression tasks.  The basic approach of PI models involves estimating a pair of quantile functions (\cite{koenker1978regression}), say ($F_{q}(x),F_{1+q-\alpha}(x)$), of the  conditional distribution $Y|X$, for some $0 \leq q \leq \alpha$, where the $q^{th}$ quantile function for given $x$ is defined as infimum of functions satisfying $P(y\leq F_q(x)|x) = q$.  

For time-series data, estimating the Prediction Interval (PI) for future observations using an auto-regressive approach is referred to as probabilistic forecasting.
 Both PI estimation and probabilistic forecasting models are widely investigated in the Neural Network (NN) architectures in the literature. The PI estimation methods in the NN literature can primarily be divided into two main categories. A popular class of PI estimation methods assumes that the conditional distribution $Y|X$ follows a particular distribution (often normal) and obtains the quantile functions by computing the inverse of the  corresponding cumulative distribution function. Some important of them are Bayesian method (\cite{mackay1992evidence, bishop1995neural}, Delta method\cite{de1998prediction,hwang1997prediction,seber2015nonlinear}) and Mean Variation Estimation (MVE) method (\cite{nix1994estimating}). Some of the recent NN architecture for the probabilistic forecasting task with distribution assumptions are Mix Density Network (\cite{bishop1994mixture,zhang2020improved}), Deep Auto-regressive Network (Deep AR) (\cite{salinas2020deepar}).  


The other class of PI estimation and probabilistic forecasting methods believe in estimating the  pair of quantile functions ($F_{q}(x),F_{1+q-\alpha}(x)$) in distribution-free setting without imposing any assumption regarding the distribution of $Y|X$. For estimation of the  $q^{th}$ quantile function, $F_q(x)$, most of them minimizes the  pinball loss function.  The pinball loss-based NN model, also known as Quantile Regression Neural Network (QRNN) (\cite{taylor2000quantile,cannon2011quantile}) is the main PI estimation method, which has been utilized in various engineering applications. The pinball loss-based NN model has been frequently applied to probabilistic forecasting of wind (\cite{wan2016direct}),  electric load (\cite{zhang2018improved,zhang2020improving}), electric consumption (\cite{he2019electricity}), flood risks (\cite{pasche2024neural}) and solar energy (\cite{lauret2017probabilistic}). 
Some of the distribution-free PI estimation  NN methods consider the minimization of a particularly designed loss function for the direct and simultaneous estimation of the bounds of the PI. Some of them are Lower Upper Bound Estimation (LUBE) NN, Quality-Driven (QD) Loss NN (\cite{pearce2018high}) and Tube loss NN (\cite{anand2024tube}).

Despite the remarkable success of neural architectures, researchers still prefer SVMs for their predictive accuracy in regression and classification tasks, particularly when dealing with small size tabular dataset. This is because SVMs explicitly incorporate regularization, guarantee a global optimal solution, and produce sparse solutions, which remain missing in the NN learning. One notable advantage of the SVM models over the NNs is that they obtain the global optimal solution by minimizing a convex function which makes their solution invariant to the initialization. However, in contrast to NN literature, there are only a few SVM methods which target the PI estimation and probabilistic forecasting tasks in the literature. We summarize the contribution of our work as follows. 
\begin{enumerate}
    \item [(a)] First, we carefully review the existing literature on PI estimation and probabilistic forecasting methods in SVM. We outline the desirable properties of an ideal PI model and compare the PI estimation and probabilistic forecasting methods in SVM against them.  After a thorough investigation, we find that only two of SVM PI methods attain the global optimal solution but, none of them achieve a sparse solution vector.  The global optimal and sparse solution are  two  main elegant properties of the  traditional SVM methods developed for classification and regression tasks. Sparse models offer improved generalization, reducing the complexity of both learning and prediction processes.

    \item[(b)] Building on this motivation, we propose a sparse SVM method for Prediction Interval (PI) estimation and probabilistic forecasting. For a given quantile $q \in (0,1)$ and target calibration level $1-\alpha$, the method computes two sparse quantile estimates, $(\hat{F}_{q}(x), \hat{F}_{1+q-\alpha}(x))$, which serve as the bounds of the PI. To achieve sparsity in the solution for $\hat{F}_{q}(x)$ and $\hat{F}_{1+q-\alpha}(x)$, the method optimizes the regularization of the norm $L_1$ alongside the pinball loss function parameterized by $q$ and $(1+q-\alpha)$ in two separate problems, each of them efficiently formulated as a linear programming problem. Consequently, the proposed approach solves a pair of linear programming problems to derive the sparse SVM solution for PI estimation and probabilistic forecasting tasks. Our Sparse SVM model enhances the PI estimation process by reducing the overall complexity of learning and prediction while preserving the classical properties of SVM, achieving both a globally optimal and sparse solution.
    
    \item [(c)] We develop a simple yet effective feature selection algorithm for PI estimation using our sparse SVM PI model. We show that our algorithm does not only successfully discard a significant percentage of features but, also improves the quality of the PI  while learning the PI for high-dimensional data. To the best of our knowledge, our algorithm is the first to address the  feature selection problem in PI estimation for tabular data, paving the way for more efficient and interpretable UQ models for quantifying the uncertainty of prediction with high dimensional data.   
    
    \item[(d)] We conduct extensive experiments on artificial, real-world benchmark datasets to empirically analyze the PI quality obtained by the both existing and newly developed PI estimation models in SVM.For high-dimensional datasets, we reveal the effectiveness of the Sparse SVM-based prediction interval (PI) model by performing feature selection using  our sparse SVM PI feature selection algorithm. Furthermore, we compare the performance of SVM-based PI methods with various deep learning models designed for probabilistic forecasting on benchmark datasets. The numerical results indicate that SVM-based approaches can achieve  similar PI quality to that of more complex deep learning models. 
    
    %superiority of the proposed Sparse SVM model over existing SVM models for the PI estimation task. Additionally, we applied the Sparse SVM model to obtain the probabilistic forecast of wind speed  and showed that it achieves superior PI quality compared to other SVM models used for probabilistic forecasting.
\end{enumerate}

 
The remainder of this paper is structured as follows. Section 2 provides a systematic review of the preliminaries concepts required for understanding of the SVM models developed for PI estimation and probabilistic forecasting.  Section-3 provides a detail description of the several SVM models for PI estimation and probabilistic forecasting,   highlighting their advantages and limitations. In Section 4, we introduce the proposed Sparse SVM models for PI estimation and probabilistic forecasting. Section 5 presents the numerical results from extensive experiments, demonstrating the effectiveness of the proposed SVM model for PI estimation and probabilistic forecasting tasks. Section 6 concludes our paper. 


\section{Related key concepts}
In this section, we will outline key concepts and methods relevant to PI estimation techniques in SVM.  
\subsection{Quantile Regression and SVM}
\begin{figure}[h]
    \centering
    \includegraphics[width=0.4\linewidth]{pin_loss.png}
    \caption{Pinball loss function}
    \label{fig:enter-label}
\end{figure}
In distribution free setting, for a given quantile $q \in (0,1)$, the quantile value is estimated by minimizing the pinball loss function, which is given by 
\begin{equation}
     \rho_q(u)  = \begin{cases}
        qu, \mbox{ ~~~if~~} u \geq 0, \\
        (q-1)u, \mbox{Otherwise.}, 
    \end{cases} 
\end{equation}
For the estimation of the conditional quantile function,  
 $u$ represents the error obtained by the subtracting the estimates $f(x_i)$ from its target values $y_i$.  For a given quantile $q \in (0,1)$, training set $T=\{(x_i,y_i): x_i \in \mathbf{R}^n, y_i \in \mathbf{R}, i =1,2,...m. \}$ and class of function $F$, let us suppose that $f_T$ is the solution of the problem $\min \limits_{f\ \in F} \sum \limits_{i=1}^{m}\rho_q(y-f(x_i))$. Tekeuchi et al., have shown that the fraction of $y$ values lying below the function $f_T(x)$ is bounded from above by $qm$ and asymptotically equals $qm$ with probability $1$ under a very general condition in their work (\cite{takeuchi2006nonparametric}).




 Given the training set, SVM models estimate the function in the form of $f(x) = w^T\phi(x) + b$, where $\phi$ maps the input variable $x$ into the high dimensional feature space, such that for any pair of $x_i$ and $x_j$ in   $\mathbf{R}^n$,  $\phi(x_i)^T\phi(x_j)$ can be obtained by the well defined kernel function $k(x_i,x_j)$. By the use of the kernel trick and representer theorem \cite{scholkopf2001generalized},  the SVM estimate $f(x) = w^T\phi(x) + b$ can be represented by the kernel generated function in the form of $\sum_{i=1}^{m}k(x_i,x)u_i + b$, where $k$ is positive-semi definite kernel \cite{mercer1909xvi}. This representation eliminates the need for explicit knowledge of the mapping $\phi$.

 \subsection{Support Vector Quantile Regression model}
The Support Vector Quantile Regression (SVQR) model minimizes the $L_2$-norm of the regularization along with the  empirical risk computed by the pinball loss function.  For $q^{th}$ quantile function estimation, it seeks the solution of the problem
\begin{equation}
    \min \limits_{(w,b)} \frac{1}{2} w^Tw + C\sum \limits_{i=1}^{m} \rho_q(y_i- (w^T\phi(x_i) +b)),
\end{equation}
which can be equivalently converted to the following Quadratic Programming Problem (QPP)
 \begin{eqnarray}
     \min \limits_{(w,b,\xi,\xi^*)} \frac{1}{2} w^Tw + C\sum \limits_{i=1}^{m} (q\xi_i + (1-q)\xi_i^{*}) \nonumber \\
     & \hspace{-110mm}\mbox{subject to,} \nonumber \\
     & \hspace{-80mm} y_i- (w^T\phi(x_i) +b) \leq \xi_i, \nonumber  \\
     & \hspace{-80mm} (w^T\phi(x_i) +b)-y_i \leq \xi_i^{*}, \nonumber \\
     & \hspace{-80mm} \xi_i, \xi_i^{*} \geq 0,~ i =~1,2,..m,
     \label{eq1}
 \end{eqnarray}
where $C \geq 0$ is the user defined parameter for trading-off the empirical risk against the model complexity. 

To efficiently solve QPP (\ref{eq1}), we often focus on obtaining the solution to its corresponding Wolfe dual problem, which is given by 
\begin{eqnarray}
     \min \limits_{(\alpha ,\beta)} \sum_{i=1}^{m} \sum_{j=1}^{m}(\alpha_i-\beta_i)k(x_i,x_j)(\alpha_j-\beta_j) - \sum_{i=1}^{m}(\alpha_i-\beta_i)y_i \nonumber \\
     & \hspace{-180mm}\mbox{subject to,} \nonumber \\
     & \hspace{-140mm}  \sum \limits_{i=1}^{m}(\alpha_i-\beta_i) = 0, \nonumber  \\
     & \hspace{-120mm}  0 \leq \alpha _i \leq Cq, ~ i =~1,2,..m, \nonumber \\
     & \hspace{-110mm} 0 \leq \beta _i \leq C(1-q),~ i =~1,2,..m,
     \label{eq2}
 \end{eqnarray}
where $(\alpha_i, \beta_i), i= ~1,2,..,m$, are Lagrangian multipliers.

After obtaining the optimal solution of the dual problem (\ref{eq2}), $(\alpha_i^{*}, \beta_i^{*}), i= ~1,2,..,m$, the $q^{th}$ quantile function is estimated by 
\begin{equation}
    f_q(x) = \sum_{i=1}^{m}(\alpha_i^{*}- \beta_i^{*})k(x_i,x) + b . \label{out1}
\end{equation}
The estimation of the bias term $b$ can be obtained by using the  KKT conditions of the primal problem (\ref{eq1}). For this, we need to select the every training point $(x_k,y_k)$ which corresponds to $0 < \alpha_k^{*} < Cq$ or $0 < \beta_k^{*} < C(1-q)$ and compute
\begin{equation}
    b_k= y_k - \sum \limits_{i=1}^{m}(\alpha_i^{*}-\beta_i^{*}) k(x_i,x_k).
\end{equation}
The final value of bias $b$ can be obtained by computing the the mean of  all $b_k$. 

\subsection{Least Squares Support Vector Regression }
For estimating the mean regression using training set $T$, the LS-SVR model \cite{suykens2002weighted} minimizes the least square loss function along with the $L_2$-norm of regularization in the following problem.
  \begin{eqnarray}
     \min \limits_{(w,b,\xi)} \frac{1}{2} w^Tw + C\sum \limits_{i=1}^{m} (\xi^2_i) \nonumber \\
     & \hspace{-90mm}\mbox{subject to,} \nonumber \\
     & \hspace{-40mm} y_i- (w^T\phi(x) +b) = \xi_i,~ i =~1,2,..m.
     \label{eq4}
 \end{eqnarray}   
The solution of problem (\ref{eq4}) can be obtained through the  following system of equations
  \begin{equation}
  \begin{bmatrix}
  0 & e^T\\
  e & K(A,A^T) + \frac{2}{C}I
  \end{bmatrix} \begin{bmatrix}
  b \\
  \alpha
  \end{bmatrix} =
  \begin{bmatrix}
  0 \\
  Y
  \end{bmatrix}, \label{eq3}
  \end{equation}
where \( K(A, A^T) \) is an \( m \times m \) kernel matrix constructed from the training set \( T \), \( e \) is an \( m \)-dimensional column vector of ones, and \( I \) represents the \( m \times m \) identity matrix.
After obtaining the $(b,\alpha)$ from (\ref{eq3}), the LS-SVR estimates the regression function for a given $x \in \mathbb{R}^n$ using
  \begin{equation}
  {f}(x) = \sum_{i=1}^{m}k(x_i,x)\alpha_i + b . \label{eq5}
  \end{equation}
  
\subsection{Probabilistic Forecasting} 
The task of probabilistic forecasting is basically an extension of the PI estimation in an auto-regressive setting. Consider the time series observations $  T = \{ x_1,x_2,....,x_t \}$, recorded on $t$ different time stamps.  If \( p < t \) denotes the effective lag window, then auto-regressive models estimate the relationship between \( z_i := (x_{i-p+1}, \ldots, x_i) \) and \( x_{i+1} \) for \( i = p, p+1, \ldots, t-1 \) using the training set \( T \), and use this learned relationship to forecast future observations. Point forecasting models aim to estimate the conditional expectation $E(x_{i+1} \mid z_i)$. However, in many cases, such forecasts may incur significant errors due to inherent noise and volatility in the data. Probabilistic forecasting quantifies these uncertainties in the prediction by obtaining the PI. 

 The task of probabilistic forecasting is to estimate the PI for $x_{i+1}$ given the input $z_{i}$ for $i \geq  t$. The SVM based probabilistic forecasting models obtain the estimate of  the PI $[\hat{F}_{q}(z_{i}),\hat{F}_{1+q-\alpha}(z_{i})]$, where $\hat{F}_{q}(z_{i})$ and $\hat{F}_{1+q-\alpha}(z_{i})$ are  kernel generated functions, estimating of the  $q^{th}$ and ${(1-\alpha-q)}^{th}$ quantiles of the  conditional distribution $(x_{i+1} \mid z_{i})$ for some $ q \in (0,\alpha)$. Distribution-free probabilistic forecasting methods estimate quantile functions directly, without making any assumptions about the conditional distribution $(x_{i+1} \mid z_{i})$.




\section{PI estimation in SVM}
In this section, we gather in detail the PI estimation and probabilistic forecasting methods developed in the SVM literature and compare their advantages and limitations. 

\subsection{PI estimation through LS-SVR}
One of the naive PI estimation method in SVM literature follows the normal assumption regarding the distribution of $Y|X$ and estimate its mean through (\ref{eq5}) by training the LS-SVR model. The error distribution \(\epsilon_i = y_i - f(x_i)\) follows a normal distribution with a mean of zero and variance \(\sigma\). This variance can be estimated from the error computed on training set $T$. The pair of quantile bounds required for PI is estimated as $(\hat{f(x)}+\epsilon_{\frac{\alpha}{2}}, \hat{f(x)}+\epsilon_{1-\frac{\alpha}{2}}) $, where $\epsilon_q$ is the $q^{th}$ quantile of the error.
A more refined and bias-corrected PI based on the LS-SVR model is proposed in (\cite{de2010approximate,cheng2014confidence}).

\subsection{PI estimation through SVQR}
Given the high confidence $1-\alpha$ with training set $T$, the PI model requires the estimation of the pair of quantile functions ($F_{q}(x),F_{1+q-\alpha}(x)$) of the  conditional distribution $Y/X$. The SVQR model can be trained twice for the estimation of the  pair of these quantile functions for some $0 \leq q \leq \alpha$.
We detail the algorithm for PI estimation through SVQR in Algorithm 1. At Algorithm 1, the tuning of $C$ refers to selecting the value of $C$ from a specified range such that the SVQR estimate obtains the least coverage error.

\begin{algorithm}
\caption{PI estimation through SVQR}\label{alg:euclid}
\begin{algorithmic}[1]
\Procedure{:- PI through SVQR} {$T,1-\alpha$}
\State Choose some $\bar{q} \in (0,1-\alpha)$ 
\For{ each $q \in \{\bar{q}, (1+\bar{q}-\alpha)\} $} 
\State Solve the QPP problem (\ref{eq2}) by tuning the value of $C$. Obtain the solution $(\alpha^*,\beta^*)$.
\State  Estimate the function $f_q(x)$ using the (\ref{out1}). 
\EndFor\label{euclidendwhile}
\State \textbf{return} $(f_{\bar{q}}(x),f_{1+\bar{q}-\alpha}(x))$
\EndProcedure
\end{algorithmic}
\end{algorithm}

A key challenge in estimating prediction intervals (PI) using the quantile approach is to determine the good choice of $\bar{q}$  for obtaining the narrower PI. For a symmetric noise distribution, \(\bar{q} = \frac{\alpha}{2}\) is expected to produce the PI with minimum width. However, this does not hold for an asymmetric noise distribution. In the latter case, \(\bar{q}\) should be selected such that the resulting PI passes through the denser regions of the data cloud. Furthermore, for each choice of $\bar{q}$, the optimization problem (\ref{eq2}) must be solved twice for $q = \bar{q}$ and  $q = 1 + \bar{q} - \alpha$ to obtain the prediction interval (PI). It increases the overall computational complexity of the PI estimation process, making it both time-consuming and challenging in practice.

To simplify the PI estimation process, researchers have developed direct PI estimation methods, which solve a single optimization problem to  obtain the both bounds of PI simultaneously. These methods are designed with a specialized loss function that can be minimized to obtain the both bounds of PI simultaneously.  Some important of them are LUBE loss (\cite{khosravi2010lower}), Quality-Driven (QD) loss (\cite{pearce2018high}) and Tube loss (\cite{anand2024tube}) functions. We describe those also formulated within the SVM framework as follows.

\subsection{PI estimation through Tube loss}
(\cite{anand2024tube}) have developed the Tube loss for PI estimation and probabilistic forecasting. It can be minimized directly to obtain the bounds of the PI simultaneously. The minimizer of the Tube loss function also guarantees the target coverage $1-\alpha$ asymptotically.  Also, the PI tube can also be shifted up or down by tuning its parameter $r$ so that it can cross through the denser region of data cloud for minimal PI width. Furthermore, the width of the PI tube can be explicitly minimized in its optimization problem through the parameter $\delta$. It helps to improve the PI width,  when the PI tube achieves a coverage higher than the target on the validation set.

 The Tube loss function is a kind of two-dimensional extension of the pinball loss function (\cite{koenker1978regression}). For a given $1-\alpha \in (0,1)$ and $u_2 \leq u_1$,  the Tube loss function is given by
 \begin{eqnarray}
		  \rho_{1-\alpha}^{r}(u_2, u_1) =
		 \begin{cases}
                 (1-\alpha) u_2,  ~~\mbox{if}~~~ u_2  > 0,\\
			-\alpha u_2,  \mbox{~~~~~~~if}~u_2 \leq 0 ,u_1 \geq 0   \mbox{~and~~}  {ru_2+(1-r)u_1} \geq 0,\\
			\alpha u_1, \mbox{~~~~~~~~~if}~u_2 \leq 0 ,u_1 \geq 0  \mbox{~and~~}  {ru_2+(1-r)u_1} < 0,\\
			-(1-\alpha) u_1, ~~ \mbox{if}~~ u_1  < 0, 
		\end{cases}
		\label{ppl1}
	\end{eqnarray}
 where $0 <r < 1$ is a user-defined parameter and ($u_2$, $u_1$) are errors, representing the deviations of $y$ values from the bounds of PI.  
 \begin{figure}[htp]
		\centering 
 {\includegraphics[width=0.50\linewidth, height = 0.30\linewidth ]{tube_09 (1) (1).png}}   
  \caption{Tube loss function for $1-\alpha$ = 0.9.}
  \label{tubeloss}
  \end{figure}
 
 Figure \ref{tubeloss} illustrates the Tube loss for $(1-\alpha) = 0.9$ with $r = 0.5$.  For $r = 0.5$ , the Tube loss function is always a continuous loss function of $u_1$ and $u_2$, symmetrically positioned around the line $u_1 + u_2 = 0$.  In all experiments with a symmetric noise distribution, the \( r \) parameter in the Tube loss function should be set to \( 0.5 \) to capture the denser region of \( y \) values. 

 The Tube loss SVM model seeks a pair of kernel generated functions
 \begin{equation}
     \mu_1(x) = \sum_{i=1}^{m}k(x_i,x)\alpha_i + b_1 \mbox{~~and~}  \mu_2(x) = \sum_{i=1}^{m}k(x_i,x)\beta_i + b_2  \label{11}
 \end{equation}
  by minimizing the optimization problem

  \begin{eqnarray}
\min_{(\alpha, \beta,b_1,b_2)} J{(\alpha, \beta,b_1,b_2)} =  \frac{\lambda}{2}(\alpha^T\alpha + \beta^T\beta) + \sum_{i=1}^{m}\rho_{1-\alpha}^{r} \big (y_i,\big(K(A^T,x_i)\alpha + b_1\big),\big( K(A^T,x_i)\beta + b_2\big)~\big)   \nonumber \\ & \hspace{-180mm} + \delta \sum \limits_{i=1}^{m}  \big|(K(A^T,x_i)(\alpha-\beta) + (b_1-b_2) \big|, \label{prob6}
	\end{eqnarray}
where $\delta, r$ and $\lambda$ are user-defined parameters and $A$ is the $m \times n$ data matrix containing the training set. Further, details on the Tube SVM problem and its minimization using gradient descent method can be found in (\cite{anand2024tube}).  
\subsection{PI estimation through LUBE loss}
The LUBE method (\cite{khosravi2010lower}) was originally developed in the NN framework but, was extended in the SVM framework later in (\cite{shrivastava2014prediction,shrivastava2015prediction}) for probabilistic forecasting of  electric price.  For the given target confidence $(1-\alpha)$, and training set $T$, the LUBE 
 SVM model seeks a pair of kernel generated functions of \ref{11}, $\mu_1(x)$ and $\mu_2(x)$, which are obtained by minimizing the following loss function
 \begin{equation}
		CWC = \frac{1}{R}MPIW \big(1+ \gamma~ (PICP) e^{-\eta(PICP  -(1-\alpha)) } \big). \label{lubecost}
	\end{equation}

 Here, MPIW is the average width of estimated PI on training set, computed by $\frac{1}{m}\sum_{i=1}^{m} (\mu_2(x_i)-\mu_1(x_i))$.  As discussed earlier, the PICP is the coverage of the estimated PI and computed using the  $PICP = \frac{1}{m}\sum_{i=1}^{m}k_i$, where
 \begin{equation}
     k_i = \begin{cases}
			1, ~\mbox{~if~~} y_i \in  [\mu_1(x) , \mu_2(x)]. \\
			0, ~\mbox{Otherwise.}
		\end{cases} \nonumber
 \end{equation}
 Further, the $\gamma(PICP) = \begin{cases}
		0, \mbox{~if ~} PICP \geq 1-\alpha, \\
		1, \mbox{~otherwise,}
	\end{cases}$ , $R$ is the range of response values $y_i$ and $\eta$ is the user-defined parameter. 
	
The major problem with the LUBE cost function (\ref{lubecost}) is that it is very difficult to be optimized because, the PICP is a step function. Khosravi et al. have solved the  LUBE cost function (\ref{lubecost}) using the Particle Swarm Optimization (PSO) (\cite{kennedy1995particle}) to estimate the PI. However, due to sub-optimal solution, high-quality PI is not always observed. Pearce et al. refine the LUBE cost function and use a sigmoidal function to approximate the PICP, allowing the application of the gradient descent method for training  the NN for PI estimating in their work (\cite{pearce2018high}).


\subsection{Comparison of PI estimation SVM models}
In Table \ref{des_prop1}, we  visualize the desirable properties for a PI estimation model and compare the SVM methods in light of  them with a detailed discussion  as follows. 


\begin{table}[]
\centering
 \begin{tabular}{|c|c|c|c|c|c|}
 \hline
  &  LS-SVR PI &  SVQR PI &  LUBE  & Tube loss    \\ \hline
   Distribution-free method  & No  & Yes  & Yes  & Yes \\
      Asymptotic guarantees  & \footnotesize {Only with normal noise}  & Yes  & No  &  Yes \\
   Direct PI estimation   &  Yes & No & Yes & Yes  \\
    PI tube movement   & No  & Yes & No & Yes   \\
     Global optimal solution   & Yes  &  Yes & No  & No  \\
      Re-calibration   &  No  & No & Yes  & Yes  \\
       Sparsity   &  No  & No & No  & No  \\
    \hline  
 \end{tabular}
 \caption{ Comparisons of the Quantile , LUBE, QD loss and Tube loss based PI estimation models }
   \label{des_prop1}
 \end{table}

\begin{enumerate}
    \item [(a)] Distribution-free method : - The LS-SVR PI model assumes that the underlying noise distribution of the data is normal and may obtain poor estimate otherwise. In the literature, distribution-based PI methods often struggle to achieve consistent performance across various datasets. The SVQR, LUBE, and Tube loss methods provide PI estimation without assuming any specific distribution, allowing them to generate high-quality PI even in the presence of non-normal noise.
     \item[(b)] Asymptotic guarantees:-  A fundamental requirement in PI estimation models is that the obtained PI should guarantee the target coverage $1-\alpha$ at least asymptomatically, which remain missing in the LUBE method. The LS-SVR PI model provides this guarantee only in presence of normal noise. The SVQR and Tube loss based PI methods provide this asymptotic guarantee.
     \item[(c)] Direct PI estimation:- As detailed in Algorithm 1, the SVQR PI model obtains the two bounds of PI by solving  pair of SVQR problems one by one. The LUBE and Tube loss-based PI models simultaneously obtain the bounds of the PI by solving a single optimization problem.
     \item[(d)] PI tube movement:- The PI tube movement is one another important desirable feature for PI estimation. This movement allows the PI to pass through the denser regions of the data cloud, helping to minimize the width of the PI, while achieving the target coverage. The centered PI is ideal only in the presence of the symmetric noise. However, in the presence of asymmetric noise in the data, the width of the prediction interval (PI) can be reduced by shifting it upward or downward without compromising its coverage.
     The SVQR and Tube loss-based PI models enable PI movement through their parameter $\hat{q}$ and $r$ respectively.

     \item [(e)] Global Optimal Solution:- One of the attractive feature of the standard SVM methods is that they guarantee a global optimal solution by solving a convex optimization problem. However, in PI estimation, only the SVQR and LS-SVR based PI model maintains this guarantee by minimizing the convex loss function in its optimization problems. The Tube loss problem (\ref{prob6}) is non-convex and hence fails to guarantee the global optimal solution. Furthermore, the LUBE problem (\ref{lubecost}) is highly discontinuous and relies on meta-heuristic algorithms for its solution. It  often makes the LUBE solution suboptimal, resulting in poor PI quality.

     \item [(f)] Re-calibration:-  As detailed in Algorithm 1, the SVQR PI model obtains the bound of PI by solving the pair of SVQR problems one-by one. It can not explicitly minimize the width of PI in its optimization problem. This limitation prevents SVQR PI from using the recalibration feature. In recalibration, PI models are retrained to reduce interval width, when  empirical coverage obtained on validation set exceeds the target $1-\alpha$ significantly. During retraining, the PI models increases the value of the parameter ($\delta$ in case of the Tube loss) that trade-off the width of the PI against the coverage in the optimization problem. The LUBE and Tube loss-based models explicitly incorporate the minimization of prediction interval (PI) width in their problems, thereby enabling recalibration, which is practically useful for further reducing the PI width.
     \item [(g)] Sparsity:- Sparsity is yet another promising feature offered by the initial SVM models developed for the classification and regression. However, it remains missing with all PI estimation models developed in the SVM framework.
\end{enumerate}
\section{Sparse  PI estimation in SVM}		
In this section, we introduce the Sparse Support Vector Quantile Regression (SSVQR) model and detail the algorithm for obtaining the sparse PI through it. The SSVQR PI model inherits all properties of SVQR listed in Table \ref{des_prop1} and also brings the sparse solution as well.  Additionally, it addresses the feature selection problem in PI estimation efficiently.

\subsection{Sparse Support Vector Quantile Regression model}

The SSVQR minimizes the pinball loss function with $L_1$- norm regularization.  It seeks the solution of the following problem
\begin{equation}
    \min \limits_{(w,b)} \frac{1}{2} ||w||_1 + C\sum \limits_{i=1}^{m} \rho_q(y_i- (w^T\phi(x_i) +b)), \label{ssvqr1}
\end{equation}
where $C \geq 0$ is the user defined parameter for trading of the regularization against the empirical loss. The solution to problem (\ref{ssvqr1}) is sparse, similar to LASSO regression, as the minimization of the \(L_1\)-norm regularization compels the weight coefficients to shrink to zero.

With the help of the kernel trick,  the representor theorem rewrites the estimated function $ f(x)=(w^T\phi(x) +b) $  as $\sum \limits_{j=1}^{m}k(x_j,x)u_j+ b$. It makes the SSVQR problem (\ref{ssvqr1}) equivalent to
\begin{eqnarray}
     \min \limits_{(u,b,\xi,\xi^*)} \frac{1}{2} ||u||_1 + C\sum \limits_{i=1}^{m} (q\xi_i + (1-q)\xi_i^{*}) \nonumber \\
     & \hspace{-110mm}\mbox{subject to,} \nonumber \\
     & \hspace{-70mm} y_i- \Big( \sum \limits_{j=1}^{m}k(x_j,x_i)u_j+ b \Big) \leq \xi_i, \nonumber  \\
     & \hspace{-70mm} \Big(\sum \limits_{j=1}^{m}k(x_j,x_i)u_j+ b\Big)-y_i \leq \xi_i^{*}, \nonumber \\
     & \hspace{-80mm} \xi_i, \xi_i^{*} \geq 0,~ i =~1,2,..m,
     \label{ssvqr2}
 \end{eqnarray}

 
Without loss of generality, let us consider the solution vector $u = r-p$, where $r$ and $p$ are vectors of positive numbers i,e., $r_i,p_i > 0 , i=1,2,..,m$ ,then the problem (\ref{ssvqr2}) can be expressed as
\begin{eqnarray}
     \min \limits_{(r, p,b,\xi,\xi^*)} \frac{1}{2} \sum\limits_{i=1}^{m}(r_i + p_i) + C\sum \limits_{i=1}^{m} (q\xi_i + (1-q)\xi_i^{*}) \nonumber \\
     & \hspace{-140mm}\mbox{subject to,} \nonumber \\
     & \hspace{-70mm} y_i- \Big( \sum \limits_{j=1}^{m}k(x_j,x_i)(r_j-p_j)+ b \Big) \leq \xi_i, \nonumber  \\
     & \hspace{-70mm} \Big(\sum \limits_{j=1}^{m}k(x_j,x_i)(r_j-p_j)+ b\Big)-y_i \leq \xi_i^{*}, \nonumber \\
     & \hspace{-80mm} \xi_i, \xi_i^{*},r_i,p_i\geq 0, ~ i =~1,2,..m.
     \label{ssvqr3}
 \end{eqnarray}
The above problem (\ref{ssvqr3}) is a Linear Programming Problem (LPP) with $4m$ variables, $2m$ linear constraints and $4m$ non-negative constraints, which can be efficiently solved by any LPP solver. The optimal solution  $(r^*,p^*,b^*)$ of the LPP (\ref{ssvqr3})  determines the estimate of the $q^{th}$ quantile function using

\begin{equation}
    f_q(x) = \sum_{i=1}^{m}(r_i^{*}- p_i^{*})k(x_i,x) + b . \label{out2}
\end{equation}
 
\subsection{PI estimation through SSVQR}
We detail the algorithm for PI estimation through SSVQR as follows.
\begin{algorithm}
\caption{PI estimation through SSVQR}\label{alg:euclid}
\begin{algorithmic}[1]
\Procedure{:- PI through SSVQR} {$T,1-\alpha$}
\State Choose some $\bar{q} \in (0,1-\alpha)$ 
\For{ each $q \in \{\bar{q}, (1+\bar{q}-\alpha)\} $} 
\State Solve the LPP problem (\ref{ssvqr3}) by tuning the value of $C$. Obtain the solution $(r^*,p^*,b^*)$.
\State  Estimate the function $f_q(x)$ using the (\ref{out2}). 
\EndFor\label{euclidendwhile}
\State \textbf{return} $(f_{\bar{q}}(x),f_{1+\bar{q}-\alpha}(x))$
\EndProcedure
\end{algorithmic}
\end{algorithm}
The SSVQR PI preserves the properties of the SVQR PI, including the global optimal solution, PI tube movement, asymptotic guarantees, and distribution-free estimation, while also achieving a sparse solution.


\subsection{Feature selection in PI estimation through SSVQR}
Similar to other machine learning tasks, the PI estimation in high-dimensional settings also presents several challenges. The increased dimensionality not only increase the complexity of the PI bounds but also necessitates a larger sample size to ensure the quality of the estimate PI. Therefore, an efficient feature selection method is crucial for reducing the overall complexity of  the PI estimation, particularly when dealing with high-dimensional data.

\begin{algorithm} 
\caption{Feature selection through SSVQR}\label{alg:euclid}
\begin{algorithmic}[1]
\Procedure{:- Feature selection through SSVQR} {$T,1-\alpha, \epsilon$}
\State Choose some $\bar{q} \in (0,1-\alpha)$ 
\For{ each $q \in \{\bar{q}, (1+\bar{q}-\alpha)\} $}
\State  Consider the linear kernel  $k(x_i,x_j) = x_i ^{T}x_j$ at the  LPP (\ref{ssvqr3}). 
\State  Solve the LPP (\ref{ssvqr3}) and obtain its solution $(r^*,p^*,b^*)$.
\State Obtain  the $w_q$ using  $ [x_1, x_2,...,x_m ](r^*-p^*)$. 
\EndFor \label{euclidendwhile}
\State  Compute $I_{\bar{q}} = \{ i:,| w_{\bar{q}}(i)| \leq \epsilon \}$ and  $I_{1+\bar{q}-\alpha} = \{ i:,| w_{1+\bar{q}-\alpha}(i) | \leq \epsilon \}$
\State Compute $I = I_{\bar{q}} \cap I_{1+\bar{q}-\alpha} $ and \textit{Feature Set} = $\{1,2,...,m\} - I$
\State \textbf{return}  (\textit{Feature Set})
\EndProcedure
\end{algorithmic}
\end{algorithm}
We detail the feature selection algorithm through SSVQR model for the  linear PI estimation task in Algorithm 3. Here, a linear PI refers to  the PI,  where both bounds are linear functions of the input variables ,i,e. $F_{\bar{q}}(x) = w_{\bar{q}}^{T}x + b_{\bar{q}}$  $F_{1+{\bar{q}}-\alpha}(x) = w_{1+{\bar{q}}-\alpha}^{T}x + b_{1+{\bar{q}}-\alpha}$.

The input of Algorithm 3 is the training set $T= \{ (x_i,y_i), x_i \in \mathbf{R}^n, y_i \in \mathbf{R}, i =1,2,....m \}$, specified confidence $(1-\alpha)$, and a very small number $\epsilon$. It finally returns the selected feature set.  In the next section, we  have explored several real-world datasets with numerous features and successfully perform feature selection using the SSVQR method for linear PI estimation, without compromising the quality of the estimated PI.

%In contrast to NN, there is a limited literature which considers the PI estimation task in the SVM framework. The SVM methods are still popular choice over the NN methods.  

\section{Experimental Results}
In this section, we present the numerical results to analyze the quality of the PI obtained by the different SVM models through a series of experiments on simulated/artificial and real-world  benchmark datasets. We also evaluate the effectiveness of the proposed SSVQR model and assess the quality of the PI estimated using it.  We apply our Algorithm 3 for feature selection in high-dimensional real-world benchmark datasets for the linear PI estimation task. We have also compared the SVM based probabilistic forecasting models with the deep forecasting models on several time-series benchmark datasets.

One of the key strengths of SVM machines is their ability to always obtain the global optimal solution, maintaining their relevance and applicability in modern cutting-edge technology.
As detailed in Table \ref{des_prop1}, PI estimation through SVQR and LS-SVR only obtains the global optimal solution but lacks the sparsity. In contrast, the proposed SSVQR PI can obtain the optimal global solution as well as the sparse solution. In view of this, we find the SVQR and LS-SVR model are qualified enough to be compared with the SSVQR model for the PI estimation task in SVM framework. The Tube and LUBE loss PI estimation methods available in SVM framework do not guarantee the global optimal solution and their solution may vary with the choice of initialization. However, we have considered the Tube loss and an improved version of the LUBE loss, the QD loss function (\cite{pearce2018high}) in deep forecasting models to compare their performance with the SVM based probabilistic forecasting models. 


\subsection {Evaluation Criteria and Parameter Tunning }
 Now, we describe in detail the evaluation criteria that will be used for our experiments.  In all of our experiments, our aim is to estimate the PI with a confidence level of $1-\alpha = 0.95$ that requires the estimation of the pair of quantile functions $(F_q(x),F_{q+0.95}(x))$, where $0 \leq {q} \leq 0.05 $. In case of artificial datasets, we know the noise distribution and the true quantile function can be easily computed by the inverse of the cumulative distribution function. Therefore, the quality of the quantile function can be accurately assessed by computing the RMSE between the true and estimated quantile functions. In the absence of information about the noise distribution, the quality of the quantile function can be evaluated using Coverage Probability (CP), which measures the fraction of \( y \) values falling below the estimated quantile function. For a \( q^{th} \) quantile estimate, the CP should be as close as possible to \( q \).  For evaluating the  overall quality of PI estimation, we use PICP and MPIW as assessment criteria. An effective PI method must achieve the target \( 1-\alpha \) calibration while minimizing the PI width, as measured by the MPIW value.  Furthermore, we define the Prediction Interval Coverage Error (PICE) as \( \max(0, (1-\alpha) - PICP) \) to quantify the extent to which the model falls short of the target calibration (1-$\alpha$). For comparing PI estimation models, a natural decision criterion is that the model with the lowest PICE should be considered the best. If all models successfully achieve the target calibration, the one with the minimum MPIW should be deemed the most optimal.

To estimate both bounds of the PI, we utilize the RBF kernel, defined as \( k(x_i, x_j) = e^{-q||x_i - x_j||^2} \). As detailed in Algorithm 1 and Algorithm 2, SVQR and SSVQR require solving the QPP (\ref{eq2}) and LPP (\ref{ssvqr3}) twice to obtain the quantile bounds of the PI respectively. We solve the QPPs of the SVQR PI model and the LPPs of the SSVQR model in MATLAB using the 'quadprog' and 'linprog' packages respectively. The SVQR problem(\ref{eq2}) or  SSVQR (\ref{ssvqr3}) problem requires the supply of the two user defined parameters $C$ and  RBF kernel parameter $q$ for non-linear PI estimation. We have tunned the value of the these parameters using the grid search in the search space $\{2 ^{-8},2^{-7},....,2^7,2^8 \} \times \{2 ^{-8},2^{-7},...,2^7,2^8 \} $. 

\subsection{Artificial Datasets}
First, we generate six distinct artificial datasets. In each dataset, the values of \( x_i \) are randomly sampled from a uniform distribution \( U(-5,5) \), while the corresponding \( y_i \) values are obtained by polluting  the function  
$
(1 - x_i + 2x_i^2) e^{-0.5 x_i^2}
$ 
 with different types of noise as described below.
\begin{eqnarray}
    \mbox{ \textbf{AD 1:-}} ~~y_i = (1 - x_i + 2x_i^2) e^{-0.5 x_i^2} + \xi_i , \mbox{where~~}  \xi_i \sim N(0, 0.6) \nonumber \\
    \mbox{ \textbf{AD 2:-}} ~~y_i = (1 - x_i + 2x_i^2) e^{-0.5 x_i^2} + \xi_i , ~~~~~\mbox{where~~}  \xi_i \sim \chi^2(3) \nonumber \\
    \mbox{ \textbf{AD 3:-}} ~~y_i = (1 - x_i + 2x_i^2) e^{-0.5 x_i^2} + \xi_i , \mbox{where~~}  \xi_i \sim N(0, 0.4) \nonumber \\
    \mbox{ \textbf{AD 4:-}} ~~y_i = (1 - x_i + 2x_i^2) e^{-0.5 x_i^2} + \xi_i , \mbox{where~~}  \xi_i \sim N(0, 0.8) \nonumber \\
    \mbox{ \textbf{AD 5:-}} ~~y_i = (1 - x_i + 2x_i^2) e^{-0.5 x_i^2} + \xi_i , \mbox{where~~}  \xi_i \sim U(-5, 5) \nonumber \\
     \mbox{ \textbf{AD 6:-}} ~~y_i = (1 - x_i + 2x_i^2) e^{-0.5 x_i^2} + \xi_i , \mbox{where~~}  \xi_i \sim U(-4, 4) \nonumber 
\end{eqnarray}
In case of each dataset, $2500$ data points are generated in which $1000$ data points are considered for training and rest of them are considered for testing.




% Please add the following required packages to your document preamble:
% \usepackage{multirow}
\begin{table}[]
{\fontsize{9}{9}\selectfont \begin{tabular}{|l|l|l|l|l|l|l|l|l|}
\cline{1-9}
                      & ($\bar{q},~1+\bar{q}-\alpha$)              & RMSE(Lw,Up)   & Spar (Lw,Up)     & CP (Lw,Up)     & PICP  & PICE & MPIW   & Time (s) \\ \cline{1-9} 
\multirow{5}{*}{SVQR} & (0.01, 0.96)   & (1.8021, 1.4446) & (0\%, 0\%) & (0.0110, 0.9680) & 0.957 & 0  & 2.9054 & 0.4783   \\  
                      & (0.015, 0.965) & (1.7046, 1.4653) & (0\%, 0\%) & (0.0150, 0.9720) & 0.957 & 0  & 2.8244 & 0.5485   \\  
                      & (0.02, 0.97)   & (1.6519, 1.4799) & (0\%, 0\%) & (0.0180, 0.9720) & 0.954 & 0  & 2.7855 & 0.4712   \\  
                      & (0.025, 0.975) & (1.5857, 1.5016) & (0\%, 0\%) & (0.0220, 0.9760) & 0.954 & 0  & 2.7379 & 0.4928   \\ 
                      & (0.03, 0.98)   & (1.5030, 1.5667) & (0\%, 0\%) & (0.0250, 0.9800) & 0.955 & 0  & 2.7212 & 0.6054   \\  \hline
                       \multirow{5}{*}{SSVQR}   &  (0.01, 0.96) & (1.8060, 1.4418) & (15\%, 18\%) & (0.0110, 0.9680) & 0.957 & 0 & 2.9056 & 0.6157 \\  
       & (0.015, 0.965) & (1.7051, 1.4714) & (15\%, 18\%) & (0.0150,  0.9700) & 0.955 & 0 & 2.8275 & 0.5844 \\ 
       & (0.02, 0.97) & (1.6497, 1.4861) & (15\%, 18\%) & 0.0160, 0.9720 & 0.956 & 0 & 2.7939 & 0.6169 \\ 
       & (0.025, 0.975) & (1.5517, 1.4978) & (15\%, 20\%) & (0.0250, 0.9750) & 0.95 & 0 & 2.6995 & 0.5546 \\ 
       & (0.03, 0.98) & (1.5133, 1.5604) & (16\%, 20\%) & (0.0290, 0.9800) & 0.951 & 0 & 2.7288 & 0.515 \\ \hline
\end{tabular}}
\caption{Performance of the SVQR and SSVQR PI models on AD 1 dataset.}
\label{AD1}
\end{table}
\begin{table}[]
{\fontsize{9}{9}\selectfont \begin{tabular}{|l|l|l|l|l|l|l|l|l|}
\cline{1-9}
        & ($\bar{q},~1+\bar{q}-\alpha$)               & RMSE(Lw,Up)   & Spar (Lw,Up)     & CP (Lw,Up)     & PICP  & PICE & MPIW   & Time (s) \\ \cline{1-9} 
\multirow{5}{*}{SVQR} &
        (0.01, 0.96) & (4.2653, 5.7073) & (0\%, 0\%) & (0.0100, 0.9470) & 0.937 & 0.013 & 8.3755 & 0.805 \\ 
        & (0.015, 0.965) & (4.1993, 6.0698) & (0\%, 0\%) & (0.0210, 0.9520) & 0.931 & 0.019 & 8.6865 & 0.699 \\ 
        & (0.02, 0.97) & (4.1312, 6.5751) & (0\%, 0\%) & (0.0220, 0.9610) & 0.939 & 0.011 & 9.1619 & 0.6595 \\
       & (0.025, 0.975) & (4.0525, 6.8163) & (0\%, 0\%) & (0.0340, 0.9650) & 0.931 & 0.019 & 9.3256 & 0.719 \\ 
       & (0.03, 0.98) & (4.0251, 7.4745) & (0\%, 0\%) & (0.0350, 0.9720) & 0.937 & 0.013 & 9.978 & 0.6809 \\ 
        \hline
       \multirow{5}{*}{SSVQR} &  (0.01, 0.96) & (4.1628, 5.5481) & (20\%, 15\%) & (0.0080, 0.9390) & 0.931 & 0.019 & 8.1036 & 0.502 \\ 
         & (0.015, 0.965) & (4.0206, 5.9335) & (18\%, 18\%) & (0.0100, 0.9520) & 0.942 & 0.008 & 8.346 & 0.5173 \\ 
        & (0.02, 0.97) & (4.0156, 6.3077) & (20\%, 25\%) & (0.0110, 0.9640) & 0.953 & 0 & 8.7801 & 0.5128 \\ 
        & (0.025, 0.975) & (3.9444, 7.1596) & (15\%, 20\%) & (0.0240, 0.9680) & 0.944 & 0.006 & 9.5723 & 0.4694 \\ 
        & (0.03, 0.98) & (3.9309, 7.3834) & (20\%, 20\%) & (0.0290, 0.9740) & 0.945 & 0.005 & 9.8075 & 0.5113 \\ \hline
    \end{tabular}}
    \caption{Performance of the SVQR and SSVQR PI models on AD 2 dataset.}
\label{AD2}
\end{table}

\begin{table}
    {\fontsize{9}{9}\selectfont \begin{tabular}{|l|l|l|l|l|l|l|l|l|}
\cline{1-9}
        & ($\bar{q},~1+\bar{q}-\alpha$)               & RMSE(Lw,Up)   & Spar (Lw,Up)     & CP (Lw,Up)     & PICP  & PICE & MPIW   & Time (s) \\ \cline{1-9} 
\multirow{5}{*}{SVQR} & 
        (0.01, 0.96) & (1.7894, 1.5837) & (0\%, 0\%) &( 0.0120, 0.9570) & 0.945 & 0 & 2.9 & 0.4872 \\ 
       & (0.015, 0.965) & (1.7355, 1.6297) & (0\%, 0\%) & (0.0130, 0.9590) & 0.946 & 0.004 & 2.887 & 0.5077 \\ 
        & (0.02, 0.97) & (1.5696, 1.6407) & (0\%, 0\%) & (0.0230, 0.9590) & 0.936 & 0.014 & 2.7223 & 0.5004 \\ 
      &  (0.025, 0.975) & (1.5070, 1.6927) & (0\%, 0\%) & (0.0280, 0.9680) & 0.94 & 0.01 & 2.71 & 0.4907 \\ 
       & (0.03, 0.98) & (1.4694, 1.7477) & (0\%, 0\%) & (0.0300, 0.9750) & 0.945 & 0.005 & 2.7394 & 0.4745 \\ 
        \hline
   \multirow{5}{*}{SSVQR} &      (0.01, 0.96) & 1.8411, 1.5781 & 20\%, 40\% & 0.0120, 0.9570 & 0.945 & 0.005 & 2.9462 & 0.4367 \\ 
         & (0.015, 0.965) & (1.8297, 1.5917) & (40\%, 30\%)& (0.0120, 0.9570) & 0.945 & 0.005 & 2.9458 & 0.4563 \\ 
        & (0.02, 0.97) & (1.6999, 1.6969) & (40\%, 40\%) & (0.0180, 0.9600) & 0.942 & 0.008 & 2.9026 & 0.4671 \\ 
       & (0.025, 0.975) & (1.6509, 1.7191) & (40\%, 40\%) & (0.0220, 0.9620) & 0.94 & 0.01 & 2.8817 & 0.4501 \\ 
      &  (0.03, 0.98) & (1.4865, 1.7721) & (30\% , 40\%) & (0.0290, 0.9700) & 0.941 & 0.009 & 2.7739 & 0.4717 \\ \hline
    \end{tabular}}
      \caption{Performance of the SVQR and SSVQR PI models on AD 3 dataset.}
\label{AD3}
\end{table}
\begin{table}
    {\fontsize{9}{9}\selectfont \begin{tabular}{|l|l|l|l|l|l|l|l|l|}
\cline{1-9}
        & ($\bar{q},~1+\bar{q}-\alpha$)               & RMSE(Lw,Up)   & Spar (Lw,Up)     & CP (Lw,Up)     & PICP  & PICE & MPIW   & Time (s) \\ \cline{1-9} 
\multirow{5}{*}{SVQR} & 
        (0.01, 0.96) & (2.5953, 2.1197) & (0\%, 0\% )& (0.0160, 0.9530) & 0.937 & 0.013 & 4.1549 & 0.723 \\ 
         & (0.015, 0.965) & (2.5595, 2.1292) & (0\%, 0\%) & (0.0160, 0.9530) & 0.937 & 0.013 & 4.1268 & 0.6713 \\ 
       & (0.02, 0.97) & (2.4221, 2.2193) & (0\%, 0\%) & (0.0200, 0.9610) & 0.941 & 0.009 & 4.0817 & 0.6704 \\ 
       & (0.025, 0.975) & (2.2561, 2.2555) & (0\%, 0\%) & (0.0290, 0.9630) & 0.934 & 0.016 & 3.9365 & 0.6815 \\ 
       & (0.03, 0.98) & (2.2211, 2.3699) & (0\%, 0\%) & (0.0310, 0.9670) & 0.936 & 0.014 & 4.0273 & 0.6343 \\ \hline
    \multirow{5}{*}{SSVQR}   &  (0.01, 0.96) & (2.5842, 2.1364) & (20\%, 40\%) & (0.0160, 0.9530) & 0.937 & 0.013 & 4.1577 & 0.4704 \\
       & (0.015, 0.965) & (2.5475, 2.1911) & (20\%, 20\%) & (0.0160, 0.9560) & 0.94 & 0.01 & 4.1802 & 0.4745 \\ 
       & (0.02, 0.97) & (2.4939, 2.2455) & (20\%, 20\%) & (0.0190, 0.9610) & 0.942 & 0.008 & 4.183 & 0.4818 \\ 
       & (0.025, 0.975) & (2.3617, 2.3037) & (40\%, 20\%) & (0.0240, 0.9640) & 0.94 & 0.01 & 4.1002 & 0.4662 \\ 
       & (0.03, 0.98) & (2.2689, 2.4147) & (40\%, 20\%) & (0.0290, 0.9700) & 0.941 & 0.009 & 4.121 & 0.4884 \\ \hline
    \end{tabular}}
   \caption{Performance of the SVQR and SSVQR PI models on AD 4 dataset.}
\label{AD4}
\end{table}


\begin{table}
    {\fontsize{9}{9}\selectfont \begin{tabular}{|l|l|l|l|l|l|l|l|l|}
\cline{1-9}
        & ($\bar{q},~1+\bar{q}-\alpha$)               & RMSE(Lw,Up)   & Spar (Lw,Up)     & CP (Lw,Up)     & PICP  & PICE & MPIW   & Time (s) \\ \cline{1-9} 
\multirow{5}{*}{SVQR} & 
        (0.01, 0.96) & (4.9694, 4.4016) & (0\%, 0\% )& (0.0100, 0.9560) & 0.946 & 0.004 & 8.1044 & 0.7917 \\
     &   (0.015, 0.965) & (4.8456, 4.4539) & (0\%, 0\%) & (0.0180, 0.9600) & 0.942 & 0.008 & 8.0257 & 0.8009 \\ 
      &  (0.02, 0.97) & (4.8100, 4.4860) & (0\%, 0\%) & (0.0190, 0.9620) & 0.943 & 0.007 & 8.0223 & 0.736 \\ 
       & (0.025, 0.975) & (4.7111, 4.5143) & (0\%, 0\%) & (0.0280, 0.9640) & 0.936 & 0.014 & 7.943 & 0.6895 \\
      &  (0.03, 0.98) & (4.6921, 4.5368) & (0\%, 0\%) & (0.0290, 0.9680) & 0.939 & 0.011 & 7.9481 & 0.6191 \\ \hline
       \multirow{5}{*}{SSVQR}      & (0.01, 0.96) & (5.2658, 4.5044) & (40\%, 30\%) & (0.0050, 0.9610) & 0.956 & 0 & 8.5344 & 0.5394 \\ 
        & (0.015, 0.965) & (4.8787, 4.5073) & (35\%, 30\%) & (0.0160, 0.9610) & 0.945 & 0.005 & 8.1171 & 0.4835 \\ 
        & (0.02, 0.97) & (4.8138, 4.5249) & (30\%, 30\%) & (0.0190, 0.9650) & 0.946 & 0.004 & 8.0659 & 0.472 \\ 
       & (0.025, 0.975) & (4.7464, 4.6037) & (25\%, 30\%) & (0.0260, 0.9670) & 0.941 & 0.009 & 8.0692 & 0.5102 \\ 
        & (0.03, 0.98) & (4.6962, 4.7637) & (30\%, 40\%) & (0.0290, 0.9720) & 0.943 & 0.007 & 8.1832 & 0.5068 \\ \hline
    \end{tabular}}
   \caption{Performance of the SVQR and SSVQR PI models on AD 5 dataset.}
\label{AD5}
\end{table}

\begin{table}
    {\fontsize{9}{9}\selectfont 
    \begin{tabular}{|l|l|l|l|l|l|l|l|l|}
\cline{1-9}
        & ($\bar{q},~1+\bar{q}-\alpha$)              & RMSE(Lw,Up)   & Spar (Lw,Up)     & CP (Lw,Up)     & PICP  & PICE & MPIW   & Time (s) \\ \cline{1-9} 
\multirow{5}{*}{SVQR} & 
        (0.01, 0.96) & (6.1510, 5.5057) & (0\%, 0\%) & (0.0110, 0.9630) & 0.952 & 0 & 10.0826 & 0.6437 \\ 
      &  (0.015, 0.965) & (6.0517, 5.5584) & (0\%, 0\%) & (0.0150, 0.9670) & 0.952 & 0 & 10.031 & 0.6798 \\ 
     &   (0.02, 0.97) & (5.9777, 5.6130) & (0\%, 0\% )& (0.0170, 0.9700) & 0.953 & 0 & 10.0141 & 1.0752 \\ 
     &   (0.025, 0.975) & (5.8879, 5.6868) & (0\%, 0\%) & (0.0230, 0.9760) & 0.953 & 0 & 9.9987 & 0.9265 \\ 
     &   (0.03, 0.98) & (5.7599, 5.7156) & (0\%, 0\%) & (0.0290, 0.9780) & 0.949 & 0.001 & 9.8878 & 0.9445 \\ \hline
        
       \multirow{5}{*}{SSVQR} &  (0.01, 0.96) & (5.9304, 5.4581) & (10\%, 15\%) & (0.0110, 0.9610) & 0.95 & 0 & 9.7944 & 0.569 \\ 
      &  (0.015, 0.965) & (5.9107, 5.5686) & (10\%, 20\%) & (0.0140, 0.9700) & 0.956 & 0 & 9.9016 & 0.5656 \\ 
      &  (0.02, 0.97) & (5.7836, 5.5880) & (15\%, 25\%) & (0.0180, 0.9700) & 0.952 & 0 & 9.7807 & 0.5766 \\ 
      &  (0.025, 0.975) & (5.6486, 5.6038) & (20\%, 20\%) & (0.0260, 0.9710) & 0.945 & 0.005 & 9.6443 & 0.5423 \\ 
      &  (0.03, 0.98) & (5.6184, 5.6506) & (10\%, 15\%) & (0.0300, 0.9760) & 0.946 & 0.004 & 9.6653 & 0.5636 \\ \hline
    \end{tabular}}
    \caption{Performance of the SVQR and SSVQR PI models on AD 6 dataset.}
\label{AD6}
\end{table}
\subsection{Numerical Results, Discussion and Analysis}
We present the performance of the SVQR and SSVQR model for PI estimation task with the different values of $\bar{q}$ for each of simulated dataset in Table \ref{AD1}-\ref{AD6}.  The rightmost column of these Tables list the different pairs of target quantiles ($\bar{q},~1+\bar{q}-\alpha$), required for the PI estimation with $0.95$ confidence level. As detailed in Section 5.1 of this paper, for artificial datasets, the quality of the estimated upper and lower quantiles of the PI can be best evaluated by computing the RMSE between the estimated quantiles and their corresponding true quantile functions. The third column of the Tables \ref{AD1}-\ref{AD6} list these RMSE for different values of $\bar{q}$ and mean of them are plotted at Figure \ref{steady_state} for lower and upper quantile estimation separately for different simulated datasets. The quantile bounds estimated by the SSVQR model are comparable to, or slightly better than, those obtained from the SVQR model. Replacing $L_2$ regularization with $L_1$ regularization in the SVM quantile regression model does not result in significantly different estimates. However, the major advantage of using the SSVQR model is its ability to obtain the sparse solution vector. Figure \ref{fig:enter-label}(a) plots the sparsity of the solution vector corresponding to the upper and lower quantile bounds obtained by the SSVQR model.  It highlights that the SSVQR model effectively reduces significant coefficients of the solution vector to near zero which enables the feature selection task in PI estimation and also simplify the overall prediction process.
\begin{figure}
     \centering
     \subfloat[]{\includegraphics[width=0.50\linewidth, height = 0.3\linewidth]{img1_lw.png}\label{<figure1>}}~~
     \subfloat[]{\includegraphics[width=0.50\linewidth, height = 0.3\linewidth]{img1_up.png}\label{<figure2>}}
     \caption{Comparison of the quality of quantile function obtained by the SVQR and SSVQR PI models.}
     \label{steady_state}
\end{figure}


\begin{figure}
    \centering
   \subfloat[]{\includegraphics[width=0.52\linewidth,height = 0.3\linewidth]{SPARS.png}}
   \subfloat[]{\includegraphics[width=0.52\linewidth,height = 0.3\linewidth]{training.png}}
    \caption{(a) Qunatile functions estimated by the proposed SSVQR model is sparse while SVQR model fails to obtain the sparse solution. (b) Average training time comparison of the SVQR and SSVQR models for PI estimation.}
    \label{fig:enter-label}
\end{figure}

We compare the overall average training time (in seconds) taken by the SVQR PI and SSVQR PI models for the PI estimation task across different artificial datasets in Figure \ref{fig:enter-label}(b).  It shows that the SSVQR requires fewer seconds train the PI model than the SVQR model.  We have solved the QPPs of the SVQR PI model and LPPs of the SSVQR model in MATLAB with 'quadprog' and 'linprog' packages respectively. 


\begin{figure}
     \centering
     \subfloat[]{\includegraphics[width=0.58\linewidth]{qvsmpiw.png}}
     \subfloat[]{\includegraphics[width=0.51\linewidth]{qvsmpiw1.png}}
     \caption{Plot of the MPIW obtained by the SVQR and SSVQR PI models against $\bar{q}$ on  (a) AD2 and (b) AD4 dataset.}
     \label{steady_state11}
\end{figure}

Another key observation from the numerical results in Tables \ref{AD1}-\ref{AD6} is the consistent performance of the SVQR and SSVQR PI models across different datasets. In all cases, both the SVQR and SSVQR PI models manage to approximately achieve $95\%$ target coverage. We plot the MPIW values obtained by the SSVQR and SVQR PI models against different values of $\bar{q}$ on dataset AD2 and AD4 in  Figure \ref{steady_state11}. On AD2 dataset, the MPIW values of the PI obtained by both models increase with $\bar{q}$. It is  evident from Table \ref{AD2} that this increase is not related to the PICP values obtained by the SVQR and SSVQR PI models. Actually, it is caused by the nature of noise present in the AD2 dataset.  The AD2 dataset contains asymmetric noise from the ($\chi^{2}$) distribution, leading to a higher density of data points in the lower region of the input-target space, which gradually decreases as we move upward.
As $\bar{q}$ decreases, the resultant PI shifts downward, passing through denser regions of the data cloud, leading to lower MPIW values. It shows that the movement of the PI tube due to change of $\bar{q}$ values may lead to the narrower PI particularly in presence of the asymmetric noise in the data. Apart from the AD2 dataset, all other artificial datasets contain noise from symmetric distributions. In these datasets, the centered PI ($(F_{0.025}(x), F_{0.975}(x))$ is expected to achieve the high quality PI. Figure \ref{steady_state11}(b) shows that at $q = 0.025$, both the SVQR and SSVQR PI models attain the lowest MPIW values on AD4 dataset.
  
\subsection{Feature Selection through SSVQR}
Next, we apply our SSVQR model for feature selection in PI estimation with linear kernel. To demonstrate this, we use five popular real-world benchmark datasets 
namely Spambase, Student Performance, Boston Housing, UCI-secom and MADELON. We use the $80\%$ of data points for training and use rest of them for testing. We set the target calibration $1-\alpha=$ 0.95.  Following Algorithm 3, we perform feature selection using the SSVQR PI model for  $\bar{q} =0.025$ and present the results in Table \ref{tab23}.
The numerical results clearly demonstrate that the SSVQR PI feature selection (as detailed in Algorithm 3) can significantly reduce the number of features while maintaining the quality of the PI, as measured by PICP and MPIW. On average, it could reduce the $69\%$ of features on considered  datasets and reduce the complexity of the PI estimation task significantly. The training time of the PI estimating task improves significantly after dropping irrelevant features through Algorithm 3. On average, it could improve the $53\%$ improvement in training time. Also,  reducing the significant numbers of features will reduce the tunning and testing time of the  PI model. 


For high-dimensional datasets such as UCI secom and MADELON, feature selection leads to a significant improvement in the quality of the estimated PI in terms of its coverage. The SSVQR  based feature selection algorithm eliminates a significant number of irrelevant features and learn the PI in  relatively much lower dimension which results in better PI coverage with linear functions. Table \ref{tab:dropped_features} lists the features dropped by the SSVQR-PI feature selection algorithm for each dataset.


%, constructed by the pair of linear functions.  This improvement arises because, as dimensionality increases, the complexity of the underlying function required to accurately capture the quantile functions for PI estimation can grow exponentially.



%By eliminating a , the complexity of the PI estimation task is reduced, enabling high-quality PIs to be estimated even using a set of linear functions. 
\begin{table}[ht]
\centering
\renewcommand{\arraystretch}{1.5} % Increase row height
\scriptsize
\begin{tabular}{|c|c|ccc|ccc|c|}
\hline
\textbf{Dataset} & \textbf{Dimension} & \multicolumn{3}{c|}{\textbf{Before Feature Selection}} & \multicolumn{3}{c|}{\textbf{After Feature Selection}} & \textbf{\% Reduced} \\
\cline{3-9}
 &  & \textbf{PICP} & \textbf{MPIW} & \textbf{ Train Time(s)} & \textbf{PICP} & \textbf{MPIW} & \textbf{ Train Time(s)} & \textbf{Features} \\
\hline
Spambase & (4601, 58) & 0.9663 & 0.9330 & 60  & 0.9653 & 0.9340 & 52 & 46\% \\
\hline
Student Perf. & (395, 16) & 0.8861 & 6.8429 & 0.48  & 0.8861 & 6.8429 & 0.42  & 73\% \\
\hline
Boston Housing & (505, 14) & 0.9406 & 23.0647 & 0.81  & 0.9406 & 23.0647 & 0.7 & 38\% \\
\hline
uci-secom & (1568, 591) & 0.8758 & 1.6146 & 16  & 0.9204 & 1.6683 & 7 & 91\% \\
\hline
MADELON & (2000, 500) & 0.7100 & 2.0007 & 58  & 0.9375 & 2.0020 & 11  & 99\% \\
\hline
\end{tabular}
\caption{Performance comparison before and after feature selection using the SSVQR PI model}
\label{tab23}
\end{table}






\begin{table}[htbp]
\centering
\scriptsize
\begin{tabular}{|l|p{14cm}|}
\hline
\textbf{Dataset} & \textbf{Dropped Features} \\ \hline
Spambase & 2, 8, 11, 13, 16, 17, 18, 19, 20, 21, 22, 23, 25, 27, 31, 35, 37, 46, 47, 48, 49, 50, 53, 54, 55, 56 \\ \hline
Student Perf. & 1, 2, 3, 4, 5, 6, 7, 8, 9, 10, 11 \\ \hline
Boston Housing & 2, 3, 4, 7, 10 \\ \hline
UCI Secom  & 1, 4-20, 22-27, 28-41, 42-50, 52-58, 60-66, 68-70, 72, 74-87, 89, 91-110, 112-132, 134, 137-138, 140-157, 159-160, 162-187, 189-203, 205-224, 226-246, 247-296, 297-332, 333-362, 364-386, 387-400, 401-418, 420-422, 424-431, 434-435, 437-438, 440-466, 467, 469-481, 483, 489-498, 501-509, 512-520, 522-538, 540-545, 546, 548-560, 563-569, 571, 573-579, 580, 582-587 \\ \hline
MADELON & 0-89, 91-227, 229-275, 277-331, 333-444, 446-499 \\ \hline
\end{tabular}
\caption{Feature Selection by SSVQR PI}
\label{tab:dropped_features}
\end{table}

\subsection{Benchmark Datasets}

\begin{table}[htp]
    \centering
    {\fontsize{9}{10} \selectfont
    \begin{tabular}{|c|c|c|c|c|c|c|c|c|}
    \hline
      &  q & Spar (Lw, Up) & CP (Lw, Up) & PICP & PICE & MPIW & Time (s) \\ \hline
    \multirow{3}{*}{SVQR}   & (0.025, 0.925) & (0\%, 0\%) & (0.028, 0.930) & 0.90 & 0 & 28.55 & 0.1347 \\
       & (0.05, 0.95) & (0\%, 0\%) & (0.052, 0.950) & 0.90 & 0.00 & 33.06 & 0.1472 \\
     &   (0.075, 0.975) & (0\%, 0\%) & (0.088, 0.972) & 0.88 & 0.02 & 38.23 & 0.1670 \\  \hline
     \multirow{3}{*}{SSVQR}  & (0.025, 0.925) & (17\%, 16\%) & (0.027, 0.927) & 0.90 & 0  & 28.62 & 0.0657 \\ 
      &  (0.05, 0.95) & (15\%, 22\%) & (0.048, 0.947) & 0.90 & 0 & 32.77 & 0.0678 \\ 
     &   (0.075, 0.975) & (14\%, 28\%) & (0.080, 0.967) & 0.89 & 0.01 & 38.24 & 0.0681 \\ \hline
      \multirow{3}{*}{LS-SVR} &  (0.025, 0.925) & (0\%,0\% )&  & 0.91 & 0 & 31.19 & 0.0064 \\ 
     &(0.050, 0.950) & (0\%,0\%) &  & 0.91 & 0 & 30.18 & 0.0067 \\
        
    & (0.075, 0.975) & (0\%,0\%) &  &0.89 & 0.01 & 31.19 & 0.0060 \\ \hline
    \end{tabular}}
    \caption{Comparison of different SVM PI estimation methods on Boston Housing}
    \label{BS}
\end{table}



\begin{table}[htp]
    \centering
    {\fontsize{9}{10} \selectfont
    \begin{tabular}{|c|c|c|c|c|c|c|c|c|}
    \hline
      &  q & Spar (Lw, Up) & CP (Lw, Up) & PICP & PICE & MPIW & Time (s) \\ \hline
    \multirow{3}{*}{SVQR}  & (0.025, 0.925) & (0\%, 0\%) & (0.0219, 0.8997) & 0.8777 & 0.0723 & 32.5492 & 2.6141 \\ 
       & (0.05, 0.95) & (0\%, 0\%) & (0.0408, 0.9436) & 0.9028 & 0.0472 & 28.6541 & 2.8405 \\ 
       & (0.075, 0.975) & (0\%, 0\%) & (0.0690, 0.9561) & 0.8871 & 0.0629 & 32.1635 & 2.5197 \\ \hline
  \multirow{3}{*}{SSVQR}   &  (0.025, 0.925) & (12\%, 12\%) & (0.0340, 0.9029) & 0.8689 & 0.0811 & 29.4962 & 0.1858 \\ 
      &  (0.05, 0.95) & (12\%, 12\%) & (0.0583, 0.9272) & 0.8689 & 0.0811 & 27.8917 & 0.1795 \\ 
      &  (0.075, 0.975) & 12\%, 12\% & (0.0777, 0.9515) & 0.8738 & 0.0762 & 30.2547 & 0.1837 \\ \hline
    \multirow{3}{*}{LS-SVR}   &    (0.025, 0.925) &0\%,0\% & & 0.9126 & 0.0374 & 28.2429 & 0.4125 \\ 
     &   (0.050, 0.950) &(0\%,0\%)& & 0.8932 & 0.0568 & 27.3308 & 0.4631 \\
      &  (0.075, 0.975) &(0\%,0\%)&  &0.9029 & 0.0471 & 28.2429 & 0.3954 \\ \hline

    \end{tabular}}
    \caption{Comparison of different SVM PI estimation methods on Concrete dataset}
    \label{tab_para}
\end{table}


We have done experiments on the two popular benchmark datasets namely Boston Housing and Concrete  and evaluate the quality of the PI estimated by the SSVQR, SVQR and LS-SVR based PI estimation methods for different value of the $\bar{q}$ with non-linear RBF kernel.   Table \ref{BS} and \ref{tab_para} contains the numerical results obtained on the  Boston Housing and Concrete datasets respectively.  We can observe that the SSVQR, SVQR and LS-SVR PI models obtains a similar quality of the PI but SSVQR models always obtain the sparse solution vector. 

\subsection{ Probabilistic Forecasting  with SVM models}
In this section, we compare the performance of the  proposed SSVQR, SVQR and LS-SVR model for the probabilistic forecasting. We also train several recent and widely adopted deep learning models for probabilistic forecasting developed in a distribution-free setting, including Quantile-based LSTM, Tube Loss LSTM, and Quality-Driven (QD) Loss LSTM models \cite{pearce2018high} to compare their performance against the SVM-based models. The QD loss  \cite{pearce2018high} is the improved version of the LUBE model which can be used minimized with the gradient descent method in deep learning architecture. 


We consider three popular time-series datasets namely Female Births (365 $\times$ 1), Minimum Temperature (3651 $\times$ 1)  and Beer Production (464 $\times$ 1). We have used the 70$\%$ of dataset as training set and rest of them are used for testing. Out of the training set, the last $10\%$ of the observations have been used for the validation set. We present the numerical results in Table \ref{tab:performance_primary}.  All of the models were asked to obtain the probabilistic  forecast for target calibration $1-\alpha = 0.95$. 

One major observation from the Table \ref{tab:performance_primary} is that the SVM based probabilistic forecasting models obtain competitive performance compared to complex LSTM based deep forecasting models after efficient tuning of its parameters. Table \ref{tab_para} presents the tuned neural architectures of the LSTM models, along with the number of weights to be optimized for each of the probabilistic forecasting models used.
The LSTM based probabilistic forecasting models are more complex and requires the optimization of thousands of weights where as the SVQR based probabilistic models are much simpler architectures and could obtain similar quality of the PI as obtained by the LSTM based models.  Figure (\ref{fig:daily_birth}) (a) and (b) shows the forecasting of the SSVQR and SVQR model on Temperature and Beer Production datasets respectively.




\begin{table}[http]
\centering
\tiny
\resizebox{\textwidth}{!}{%
\begin{tabular}{|l| l| c| c| c| c|}
\hline
\textbf{Dataset} & \textbf{Method} & \textbf{PICP} & \textbf{MPIW} & \textbf{Training Time} & \textbf{Sparsity} \\ \hline
\multirow{5}{}{\textbf{Female Births}} & SSVQR      & 0.93  & 28.00    & 1.12    & 61 $\%$ \\ 
 & SVQR        & 0.95  & 27.11    & 0.97    & 0 $\%$ \\ 
 & LS-SVR       & 0.95  & 37.70    & 3.60    & 0 $\%$  \\ \cline{2-6}
 & Quantile LSTM & 0.95  & 28.20    & 118.00  & -- \\ 
 & Tube LSTM    & 0.96  & 28.09    & 43.00   & -- \\ 
 & QD LSTM      & 0.94  & 38.98    & -    & -- \\ \hline
\multirow{5}{}{\textbf{Minimum Temp.}} 
 & SSVQR      & 0.96  & 10.72    & 200.79  & $ 69\%$ \\ 
 & SVQR        & 0.96  & 75.81    & 172.00  & 0 $\%$ \\ 
 & LS-SVR       & 0.95  & 10.59    & 0.53    & 0 $\%$ \\ \cline{2-6}
 & Quantile LSTM & 0.95  & 24.82    & 1135.00 & -- \\ 
 & Tube LSTM    & 0.94  & 15.56    & 447.00  & -- \\ 
 & QD LSTM      & 0.79  & 5.94     & -   & -- \\ \hline
\multirow{5}{}{\textbf{Beer Prod.}} 
 & SSVQR      & 0.96  & 0.94     & 1.77    & 70 $\%$ \\ 
 & SVQR        & 0.95  & 75.73    & 0.78    & 0 $\%$ \\ 
 & LS-SVR       & 0.96  & 79.48    & 0.29    & 0$\%$  \\  \cline{2-6}
 & Quantile LSTM & 0.94  & 134.80   & 132.80  & -- \\ 
 & Tube LSTM    & 0.95  & 42.91    & 89.60   & -- \\ 
 & QD LSTM      & 0.96  & 159.71   & -    & -- \\ \hline
\end{tabular}%
}
\caption{Comparing the performance of the different SVM probabilistic forecasting models along with deep forecasting forecasting models}
\label{tab:performance_primary}
\end{table}


\begin{table}[ht]
\centering
\tiny
\resizebox{\textwidth}{!}{%
\begin{tabular}{|l|l|l|l|l|}
\hline
\textbf{Dataset} & \textbf{Model} & \textbf{Architecture} & \textbf{Weights} & \textbf{W. Size} \\ \hline
\multirow{8}{*}{\textbf{Female Births}} 
 & SVQR         & Kernel    & 256  & 10 \\ 
 & SSVQR        & Kernel    & 256  & 15 \\ 
 & LS-SVR       & Kernel    & 256  & 20 \\ \cline{2-5}

 & Quantile LSTM& LSTM [100]     & 30K  & 25 \\ 
 & Tube LSTM    & LSTM [100]     & 30K  & 25 \\ 
 & QD LSTM      & LSTM [100]     & 30K  & 25 \\ \hline
\multirow{8}{*}{\textbf{Minimum Temp.}} 
 & SVQR         & Kernel    & 2556  & 12 \\ 
 & SSVQR        & Kernel    & 2556  & 18 \\
 & LS-SVR       & Kernel    & 2556  & 22 \\ \cline{2-5}

 & Quantile LSTM& LSTM  [16,8]     & 32K  & 28 \\ 
 & Tube LSTM    & LSTM   [16,8]    & 32K  & 28 \\ 
 & QD LSTM      & LSTM   [16,8]    & 32K  & 28 \\ \hline
\multirow{8}{*}{\textbf{Beer Prod.}} 
 & SVQR         & Kernel    & 325  & 8 \\ 
 & SSVQR        & Kernel    & 325  & 14 \\ 
 & LS-SVR       & Kernel    & 325  & 18 \\ \cline{2-5}
 
 & Quantile LSTM& LSTM [64,32]      & 29K  & 24 \\ 
 & Tube LSTM    & LSTM  [64,32]      & 29K  & 24 \\ 
 & QD LSTM      & LSTM [64,32]      & 29K  & 24 \\ \hline
\end{tabular}%
}
\caption{ Comparison of the complexity of  used SVM and deep learning based probabilistic forecasting models on benchmark datasets. LSTM  [16,8] means that the used LSTM model has two hidden layer containing 16 and 8 neurons respectively. }
\label{tab:updated_comparison}
\end{table}





%\begin{table}[htbp]
%\centering
%\caption{Other Parameters for First Four Methods}
%\label{tab:performance_parameters}
%\tiny
%\resizebox{\textwidth}{!}{%
%\begin{tabular}{|l| l |c| c| c| c|}
%\hline
%\textbf{Dataset} & \textbf{Method} & \textbf{Window Size} & \textbf{S} & \textbf{c3} & \textbf{c1} \\ \hline
%\multicolumn{6}{l}{\textbf{Female Births}} \\ \hline
% & LP-SVQR      & 3    & 0.01   & 32  & 5   \\ 
% & ESVQR        & 5    & 19.00  & --  & 15  \\ 
% & LS-SVR       & 3    & 0.02   & --  & 128 \\ \hline
%\multicolumn{6}{l}{\textbf{Minimum Temperature}} \\ \hline
% & LP-SVQR      & 2    & 0.01   & 64  & 8   \\ 
% & ESVQR        & 12   & 16.00  & --  & 10  \\ 
% & LS-SVR       & 12   & 0.03   & --  & 128 \\ \hline
%\multicolumn{6}{l}{\textbf{Beer Production}} \\ \hline
% & LP-SVQR      & 3    & 0.01   & 32  & 5   \\ 
% & ESVQR        & 2    & 20.00  & --  & 13  \\ 
% & LS-SVR       & 1    & 0.02   & --  & 128 \\ \hline
%\end{tabular}%
%}
%\end{table}
\begin{figure}[htbp]
    \centering
     \subfloat[]{\includegraphics[width=1.0\textwidth, height = 0.30\textwidth]{PI in SVM/min_temp.png}} \\
     \subfloat[]{ \includegraphics[width=1.0\textwidth, height = 0.30\textwidth]{PI in SVM/beer_production.png}}
    \caption{Probabilistic forecasting with SSVQR model on daily 
 (a) Temperature and (b)  Beer Production dataset}
    \label{fig:daily_birth}
\end{figure}





\section{Conclusions}





%% The Appendices part is started with the command \appendix;
%% appendix sections are then done as normal sections
\newline
\appendix



 

\subsubsection*{Acknowledgments}
Use unnumbered third level headings for the acknowledgments. All
acknowledgments, including those to funding agencies, go at the end of the paper.
Only add this information once your submission is accepted and deanonymized. 

 \bibliography{cas-refs}
 \bibliographystyle{tmlr}

\appendix
\section{Appendix}
You may include other additional sections here.

\end{document}
